\section{Conclusions}
\paragraph{1. Conservation cannot certify accuracy.} A nearly exact invariant can coexist with negative concentrations and enormous error. Conservation, non-negativity, full-state error, and endpoint error must be reported separately, while recognising $E_{end}\le E_{global}$.

\paragraph{2. Frozen-Jacobian analysis predicts a useful scale, not an exact threshold.} The $5.895\times10^{-4}$ estimate and $5.4159\times10^{-4}$ largest stable tested EE step differ by about eight percent. This supports local linearisation as a guide, not a nonlinear proof.

\paragraph{3. Greater stability did not imply lower cost under the nominal operation-count model.} In the tested accuracy range and under the equal-weight model $C=N_f+N_J+N_{\mathrm{solve}}$, first-order IE was more expensive than EE at targets reached by both. This conclusion is restricted to the numerical methods, error targets, tested range, and nominal cost model used in this project; it is not a universal statement about implicit methods.

\paragraph{4. Tolerances can change algorithmic behaviour.} A loose Newton threshold can accept an explicit-like predictor, and an unscaled adaptive norm can allow physically invalid steps. Verified findings and proposed mixed-norm improvements must remain distinct.

Subject to final reference-table confirmation,
\[
\boxed{y_2(40)=9.185535\times10^{-6}}.
\]
